\chapter{表观遗传：序列没变，使用方式变了}

\begin{kaichang}
厨房里那瓶蜂蜜，先别急着抹面包。想一想它的来历：一个蜂巢里住着一只蜂王和几万只工蜂。蜂王有发育完全的卵巢，一天能产上千枚卵，可以活好几年；工蜂也是雌蜂，却几乎没有生育能力，夏天只能活几十天，一辈子采蜜、筑巢、照料幼虫。

关键在于：蜂王和工蜂可以是同一只蜂王生下的女儿，基因组几乎一模一样。让它们走上完全不同命运的，不是基因序列——而是幼虫期吃的东西。一直吃蜂王浆的幼虫长成蜂王；只吃两三天蜂王浆、之后改吃蜂蜜和花粉的幼虫长成工蜂。

同一套图纸，怎么盖出两栋不同的楼？更奇怪的是，楼盖好之后还能“记住”自己的样子，而在必要时又能推倒重来。这一章，我们就来拆开这套写在 DNA 之上、却不改动 DNA 的批注系统。
\end{kaichang}

\section{同一套图纸，不同的活法}

先不急着下结论，多看几个看得见、测得到的现象。

蜜蜂的例子你已经知道了。再看一对同卵双胞胎：他们从同一个受精卵分裂而来，序列几乎完全相同，小时候像得让人分不清。可几十年后，一个常年吸烟，一个坚持户外运动，两人的血压、皮肤、慢性病风险可以差得很远。去测他们的 DNA，序列还是那套序列。

再看猫。三花猫和玳瑁猫身上黑、黄、白的色块分布毫无规律，而且几乎全是母猫。更有趣的是 2001 年诞生的克隆猫“CC”：它的全部 DNA 拷贝自一只三色母猫，毛色花纹却和“原件”明显不同。序列完全相同，花纹却没有复印下来。

还有你自己。皮肤划破后愈合，长出来的是皮肤细胞，绝不会长成肝细胞或神经细胞。你体内的干细胞分裂了几十年，后代始终“记得”自己是谁。而一条毛毛虫在蛹里几乎化整为零，重组出一只蝴蝶——基因组从头到尾没变过一个字母。

类似的还有豌豆蚜：雌蚜可以不经交配生下女儿，同一批女儿在食物充足、住得宽敞时不长翅膀，一旦拥挤挨饿，后代就长出翅膀准备搬家。还是同一份基因组，按环境拨动不同的发育开关。

这些现象有一个共同点：\textbf{DNA 序列没变，基因的使用方式变了，而且这种改变能够维持}。第 70 章讲过，转录因子和染色质开放程度决定了基因此刻开不开；但当时留下了一个问题——“开放或关闭”的状态凭什么被细胞记住、一代一代传给自己的后代？答案就藏在本章的主角身上。

\section{把镜头推进到染色体上}

回忆第 66 章的图景：你的 DNA 展开有两米长，要塞进直径不到百分之一毫米的细胞核，靠的是一层层打包——DNA 像线一样缠在组蛋白“线轴”上，形成一个个核小体，核小体再折叠、盘旋成染色质。缠得紧的地方，转录机器挤不进去，基因沉默；缠得松的地方，启动子（第 68 章）暴露出来，基因才有机会被读取。

所谓表观遗传批注，主要是给这套打包系统做的两类化学记号。

\textbf{第一类：DNA 甲基化。}在胞嘧啶（C）的第 5 位碳上，由酶装上一个甲基（\chem{CH_3}），变成 5-甲基胞嘧啶。这件事几乎只发生在“C 后面跟着 G”的位置，记作 CpG。甲基从哪儿来？由一种叫 SAM 的分子快递上门——它是细胞用甲硫氨酸造出的通用“甲基供体”。注意：A、T、G、C 四个字母一个都没改，只是某些 C 的头上多了一顶化学“小帽子”。

\textbf{第二类：组蛋白修饰。}组蛋白有一截松软的“尾巴”伸在核小体外面，尾巴上的赖氨酸等氨基酸可以被加上乙酰基、甲基等小基团。乙酰化会中和赖氨酸的正电荷，削弱它与带负电的 DNA 之间的静电吸引，于是线缠得松一些；组蛋白甲基化则更微妙——同一条尾巴上，位置不同、读取它的蛋白不同，效果可以是开也可以是关。

此外还有第三类操作：染色质重塑复合体消耗 ATP 的能量，把核小体整段挪位置、换组蛋白变体。三类手段合起来，就是染色质上的“批注语言”。

\begin{table}[htbp]
\centering
\begin{tabular}{lll}
\toprule
机制 & 作用位置 & 典型效果 \\
\midrule
DNA 甲基化 & CpG 上的胞嘧啶 & 启动子区域高甲基化通常抑制转录 \\
组蛋白乙酰化 & 组蛋白尾巴的赖氨酸 & 染色质变松，倾向于促进转录 \\
组蛋白甲基化 & 组蛋白尾巴特定位置 & 可开可关，取决于位置和读取蛋白 \\
染色质重塑 & 整个核小体 & 挪动或更换核小体，改变可及性 \\
\bottomrule
\end{tabular}
\end{table}

要强调一句：这些是概率调节，不是铁律开关。某个启动子甲基化程度高，只是让它\textbf{更难}被转录，而不是物理上把门锁死；最终开不开、开多大，还要和第 70 章讲的转录因子网络叠加着算。表达量是一个连续变量，批注改变的是它的取值范围。

还有一个重要细节：人类大约六成基因的启动子附近有一段 CpG 特别密集的区域，叫 CpG 岛。健康细胞里，这些岛通常保持低甲基化，让基因处于“可用”状态；一旦某个岛被异常甲基化，对应的基因就可能长期沉默——后面讲癌症时你会再见到这个模式。反过来，基因编码区内部的甲基化却常常和活跃表达并存。可见“甲基化等于关闭”只是一句粗略口诀，批注落在什么位置，本身就是意义的一部分。

而且这套批注系统是“活”的，有一整套工具酶在维护它：有负责写入甲基的“书写者”，有把甲基逐步氧化、最终移除的“擦除者”，还有专门识别某种标记、并招募其他蛋白前来的“读取者”。三类蛋白让批注成为动态过程——细胞既能长期记住一个状态，也能在信号到来时精准改写。比如激素或发育信号可以把“擦除者”派往特定基因，清掉沉默批注，让那份基因重新进入可用名单。第 70 章讲的调控网络负责“决策”，本章讲的批注系统负责把决策“存档”。

\begin{shiyan}
\textbf{观察“同一基因组，两种颜色”}

材料：一头大蒜、两个浅盘、清水、一个不透光的纸箱。

步骤：蒜瓣剥皮，分别摆进两个浅盘，加少量水（不要没过蒜瓣）。一盘放窗台见光，一盘罩上纸箱遮光。每天补水，观察一到两周。

观察点：见光组长出绿色蒜苗，遮光组长出嫩黄的“蒜黄”——菜市场里真就卖这两种东西。两盘蒜苗的 DNA 序列一模一样，颜色差异来自叶绿素合成相关基因的表达被光打开或压住。再把蒜黄拿出来见光，几天后会转绿：这个开关是可逆的。

严格说，光的开关主要靠第 70 章讲的转录调控；但植物确实用表观批注“记”环境——冬小麦必须经历过寒冬才会抽穗开花（春化作用），细胞用组蛋白标记把某个促开花抑制基因长期关闭，这就是对寒冷的记忆。

安全提示：本实验无危险操作。只提醒一点：发霉变味的蒜不能吃，实验材料做完请丢弃。
\end{shiyan}

\section{细胞怎样把批注抄给下一代细胞}

现在回答本章最核心的问题：记忆怎样维持。

麻烦出在复制上。细胞分裂前，DNA 按半保留方式复制（第 67 章）：旧链留着，新链是现合成的——而新链上的胞嘧啶全部没有甲基。如果没人管，每分裂一次，甲基化批注就稀释一半，几代之后就荡然无存。

细胞的解决办法是一位“照抄员”：维持性甲基转移酶（DNMT1）。它专门识别“旧链有甲基、新链没有”的半甲基化 CpG，照着旧链的样子，把甲基补到新链对应的位置上。于是甲基化图案随着每次细胞分裂被誊写一遍：皮肤干细胞分裂出的后代还是皮肤细胞，几十年如一日。

\begin{tikzpicture}
\draw[thick] (0,1.6)--(4.5,1.6);
\draw[thick] (0,1.2)--(4.5,1.2);
\fill (1,1.6) circle (2pt); \fill (2.8,1.6) circle (2pt);
\fill (1,1.2) circle (2pt); \fill (2.8,1.2) circle (2pt);
\node at (2.2,2.1) {亲代：两条链都带甲基};
\draw[->,thick] (5.0,1.4)--(6.1,1.4);
\draw[thick] (6.5,2.2)--(11,2.2);
\draw[thick] (6.5,1.8)--(11,1.8);
\fill (7.5,2.2) circle (2pt); \fill (9.3,2.2) circle (2pt);
\draw (7.5,1.8) circle (2pt); \draw (9.3,1.8) circle (2pt);
\node at (8.75,2.8) {复制后：新链（空心点）暂无甲基};
\draw[thick] (6.5,0.7)--(11,0.7);
\draw[thick] (6.5,0.3)--(11,0.3);
\fill (7.5,0.7) circle (2pt); \fill (9.3,0.7) circle (2pt);
\fill (7.5,0.3) circle (2pt); \fill (9.3,0.3) circle (2pt);
\node at (8.75,-0.4) {维持酶照旧链补上：图案复原};
\end{tikzpicture}

（实心点表示甲基标记，线条表示 DNA 链——这是人为的表示法，真实的 DNA 是双螺旋，甲基是肉眼看不见的小基团。）

这套誊写系统的精度有多高？这正是理解表观遗传全部特性的钥匙。

\begin{qiaoliang}
算一笔账。人类基因组大约有 2800 万个 CpG 位点，其中大部分在大多数细胞里带着甲基。维持性誊写的保真度，实验估计大约在每位点每次分裂 96\%～99\%——听起来很高，但取 99\% 算一算：某个位点经过 50 次细胞分裂仍保持原状的概率约为 $0.99^{50}\approx 0.61$。也就是说，到你成年时，你身体里相当一部分位点的批注已经“漂移”过了（细胞里还有主动去甲基化的酶在擦除批注，所以漂移并不完全随机）。

再和第 72 章的数字对比：DNA 序列的突变率大约是每碱基每代 $10^{-8}$ 量级。表观批注的出错率比序列突变高出成千上万倍。这就是表观系统的本质：它是一块\textbf{可擦写、灵活、但不太可靠}的白板，适合记录“当前这份细胞的使用手册”，不适合当传世档案。几十亿年的长期信息，还是得刻在 DNA 序列里。
\end{qiaoliang}

还有一个壮观的例子：X 染色体失活。雌性哺乳动物有两条 X 染色体，基因剂量是雄性的两倍。早期胚胎的每个细胞会随机“封存”其中一条，压缩成几乎不表达的“巴氏小体”，并用甲基化等批注终身维持、随分裂传下去。关键在于哪个细胞封哪条是随机的——三花猫的毛色基因恰好位于 X 染色体上，封了“黑色那条”的细胞长出黄毛，反之长出黑毛，于是一只猫身上拼出随机的色块地图。这也解释了克隆猫 CC 花纹为什么不同：克隆用的供体细胞核早已做好了“封哪条”的选择，花纹模式自然和原件对不上。而三花猫几乎全是母猫，因为这套拼图需要两条 X。

\section{科学家的证据：一窝颜色各异的小鼠}

批注能改变性状，这话不能只靠逻辑推演，要有实验。下面这个故事是表观遗传学的经典证据，它的设计、排查和争论都值得细看。

\begin{zhengju}
\textbf{问题。}有一种“刺鼠”（agouti）小鼠，携带一个特殊的等位基因 $A^{vy}$。奇怪的是，同窝出生、基因型完全相同的小鼠，毛色却从纯黄到棕褐样样都有，连胖瘦都不一样：黄色个体往往肥胖、容易得糖尿病。基因型相同，性状为何不同？

\textbf{线索。}测序发现，$A^{vy}$ 的上游插着一段病毒留下的“跳跃序列”（一种转座子），它像一个不受控的开关，让毛色基因在错误的时间地点表达。再测这批序列的甲基化：棕色小鼠这里甲基化高、基因表达低；黄色小鼠甲基化低、表达高。序列一模一样，批注不同。

\textbf{实验。}2003 年，沃特兰和杰特尔给怀孕的母鼠饲料中添加甲基供体（叶酸、维生素 B12、胆碱、甜菜碱），结果子代中棕色个体的比例明显上升。也就是说，孕期饮食通过影响早期胚胎的甲基化，拨动了子代的毛色和代谢。

\textbf{替代解释排查。}会不会是甲基供体直接诱发了序列突变？测序排除了。会不会是毒理作用？所用剂量远低于毒性范围，且效果集中在特定基因座。会不会只是子宫环境的笼统影响？后续的胚胎操作实验把作用窗口锁定在早期胚胎——那正是全基因组批注大规模重写的时期。

\textbf{争论与修正。}这一类“饮食改变表观”的实验，不少样本量偏小，而且依赖转座子这类特殊序列——人类基因组里类似的构造很少。所以不能把结论外推成“孕妇吃什么决定孩子一辈子的基因开关”。这个实验真正站得住的结论是：表观批注在发育早期可以被环境拨动，而批注差异能稳定地改变性状。仅此一条，已经足够惊人。
\end{zhengju}

\section{这套系统的边界在哪里}

任何模型都有适用范围，表观遗传的边界尤其值得划清楚，因为它是被夸大最多的领域之一。

\textbf{边界一：两次“大清洗”。}哺乳动物的生命周期里，全基因组的甲基化要经历两轮大规模重置：一轮在受精后的早期胚胎，一轮在生殖细胞形成时（旧的印记被抹掉，按个体的性别重新建立）。绝大多数后天获得的批注根本活不过这两道关。这其实是好事——每一代都能从接近空白的状态重新开始，而不是背着祖辈的全部“使用痕迹”。

\textbf{边界二：例外是存在的，但都是特例。}基因组印记是一类官方许可的例外：人类大约有一百多个基因带“父源”或“母源”标签，能逃过部分重置，比如 \emph{IGF2} 基因通常只从父亲给的拷贝表达。这证明跨代携带批注在机制上可能，但它是被严格调控的少数派。植物的边界则更宽：植物的重置不如哺乳动物彻底，表观变异更容易传几代。最漂亮的例子是柳穿鱼：林奈时代就记录过一种花形突变（两侧对称的花变成辐射对称），1999 年科学家证明它不是序列突变，而是 \emph{Lcyc} 基因被甲基化沉默——它能遗传几代，偶尔又自发恢复。这是自然界真实存在的“表观突变”，但它恰恰也暴露了表观遗传的软肋：不如序列稳定，会自己翻转回去。

\textbf{边界三：人类的跨世代证据必须谨慎解释。}最著名的案例是荷兰“饥饿冬天”：1944 到 1945 年的大饥荒中，孕期挨饿的母亲生下的孩子，几十年后某些基因（如 \emph{IGF2}）的甲基化水平有微小差异，心血管代谢风险略高。但请注意两点：效应很小；而且母亲怀孕时，胎儿本人和胎儿体内的生殖细胞都被饥荒直接影响了——严格说这不叫跨世代，只能算“两代同堂的直接暴露”。要证明批注传到了从未暴露的孙辈、曾孙辈，目前的人类研究效应微弱、混杂因素太多（家庭环境、饮食传统本身也会传代，很难剥离）。诚实的结论是：真正的跨世代表观遗传在人类中\textbf{可能有，但尚无定论}。植物和低等动物里的证据更扎实，可那不能直接平移到人。

\textbf{边界四：它没有推翻任何人。}DNA 序列仍是唯一被证明能跨越千万年稳定携带信息的分子；孟德尔定律描述的等位基因分配照常成立，$A^{vy}$ 小鼠的批注再花哨，这个等位基因本身还是按孟德尔方式传下去的；自然选择作用在可遗传的表型差异上，如果某种表观变异可遗传，它同样可以被选择——只是由于保真度低，它很难在演化长跑中替代序列变异。表观遗传是基因调控系统之上的一层“可擦写内存”，是\textbf{补充}，不是革命。

\begin{wuqu}
“表观遗传证明拉马克是对的：你练出的肌肉、受过的创伤、甚至你的想法，都能写进基因传给子孙后代。”——这是关于表观遗传流传最广的过度引申。

反例与事实：第一，哺乳动物的生殖细胞要经历大规模重置，绝大多数后天批注根本到不了下一代，你的肌肉记忆只留在你自己的肌肉里。第二，人类中宣称“跨代表观遗传”的研究效应普遍微小，而且几乎无法排除更平凡的解释——创伤家庭的孩子生活在受创伤影响的环境里，环境和文化本身就会传代，不需要借用 DNA。第三，即便在证据更强的植物和低等动物中，跨代表观遗传也是概率性的、可逆的、几代之内的现象，而不是“任意改写、稳定遗传”。你的经历确实可能微调你某些细胞的批注，影响你自己的身体；但把它说成“思想可以遗传”，目前没有证据。科学最诚实的部分是：还没搞清楚的地方，就明说还没搞清楚。
\end{wuqu}

\section{回到那瓶蜂蜜}

现在可以完整回答开场的问题了。2008 年，研究者用 RNA 干扰（第 68 章提过的调控 RNA 的近亲技术）降低蜜蜂幼虫体内的 DNA 甲基转移酶，结果：本该长成工蜂的幼虫，大多数长出了发育良好的卵巢，朝着蜂王的方向发育。甲基化系统正是命运分岔的执行者之一。蜂王浆中的成分通过影响早期幼虫的表观批注（也影响激素信号），把同一套基因组拨向另一套使用方案。

再注意一个细节：工蜂的女儿如果重新被喂蜂王浆，照样能长成蜂王。批注可以改写，图纸从未动过——这正是“表观”二字的含义：在基因（epi-，之上）之上的一层调节。

双胞胎的谜也解开了：出生时两人批注高度相似，随着年龄、饮食、吸烟等差异，批注各自漂移，渐行渐远——测序会发现序列还是同一个序列，可甲基化图谱已经不一样了。

\section{可擦写，所以能入药}

表观批注“可擦写”这一点，在医学上是福音。

癌细胞有一种典型的批注错乱：全基因组甲基化大面积丢失（染色体变得不稳定），同时局部高甲基化把抑癌基因强行关闭。既然沉默是批注造成的，而不是序列丢失，就有可能逆转。药物阿扎胞苷和地西他滨掺入 DNA 后会“抓住”甲基转移酶，强迫细胞降低甲基化，让被沉默的抑癌基因重新上岗——它们已是骨髓增生异常综合征和某些白血病的标准治疗。还有一类药物抑制组蛋白去乙酰化酶，从组蛋白一侧撬动批注。

另一个热门应用是“表观遗传时钟”：2013 年，霍瓦特发现几百个 CpG 位点的甲基化水平能相当准确地预测一个人的生物学年龄，如今被广泛用于衰老研究。但要小心解读：时钟是统计相关，把时钟指针往回拨，不等于逆转了衰老本身。

印记基因出问题，还会直接导致疾病。15 号染色体上有一段区域，父源拷贝和母源拷贝各自沉默掉不同的基因。如果父源那段缺失，孩子会患普拉德-威利综合征：肌无力、食欲永远得不到满足；如果母源那段缺失，则是面貌完全不同的安格尔曼综合征：发育迟缓、频繁发笑。同一段染色体，只因为来自父亲还是母亲，后果截然不同——这是基因组印记在人类身上留下的硬证据。

克隆动物成功率低，一大原因也是供体细胞核里的批注来不及被卵子彻底重置；而把普通体细胞“洗回”干细胞的诱导多能干细胞技术，本质上就是人为洗掉细胞的表观记忆。

这带出最后一个问题：批注被重置到接近空白之后，一个受精卵怎样从零开始，用基因调控网络一步一步画出整个身体的蓝图——先分出头尾，再定出四肢，最后让三万亿个细胞各就各位？那是下一章的故事：第 75 章《从受精卵到身体》。

\begin{tiaozhan}
组蛋白尾巴上可以被修饰的位点有几十个，修饰种类有十几种，理论上的组合是个天文数字。于是有人提出“组蛋白密码”假说：这些组合构成一套密码，由特定蛋白逐个解读，决定基因命运。这个假说至今有争议。反对者指出：很多组蛋白标记是转录的“结果”而非“原因”——基因先被打开，标记才跟上；而且同一位点的标记在不同细胞类型里读法不同，相关性不等于因果。目前更稳妥的图景是：表观标记是调控网络的组成部分和记忆载体，参与反馈回路，而不是一套独立于网络之外、可以逐字翻译的密码。科学在这里还没有最终答案——这正是前沿本来的样子。跳过本节不影响主线。
\end{tiaozhan}

\section*{练一练}

\begin{sikaoti}
\item 画一张信息流图：从“幼虫吃蜂王浆”到“长成蜂王”，画出环境信号、表观批注、基因表达、发育结局之间的箭头链，并标出 DNA 序列出现在哪一步（提示：它全程不变）。
\item 用半保留复制的图景解释：为什么甲基化标记能跨细胞分裂维持，而“蛋白质本身的折叠形状”通常不能？画出复制后的旧链、新链，以及维持性甲基转移酶的作用位置。
\item 估算：若某个 CpG 位点每次细胞分裂保持原状的概率是 0.98，那么分裂 100 次后仍保持原状的概率大约是多少？（提示：$0.98^{100}$，可用 $\ln 0.98\approx -0.02$ 估算。）再和每碱基每代约 $10^{-8}$ 的序列突变率对比，说说表观批注适合承担什么任务、不适合承担什么任务。
\item 三花猫几乎都是母猫。用 X 染色体失活解释原因，并预测：如果用一只三花猫的体细胞克隆一只猫，克隆猫的毛色花纹会和它一模一样吗？为什么？
\item 一位同学读完“母亲饮食改变小鼠毛色”的新闻后说：“所以我妈怀孕时吃了什么，决定了我一辈子的基因。”请指出这句话里至少两处不严谨的地方。
\item 设计鉴别实验：你发现某种植物的性状变异能遗传，想知道它是序列突变还是表观突变。列出两个可以区分两者的检验方法，并说明各自依据的原理（提示：测序；用去甲基化药物处理后观察性状是否恢复）。
\end{sikaoti}
